% =====================================================================
%  Evaluation protocol -- proper scores for mesh-valued ensembles
%  Compile through main.tex.
% =====================================================================

\section{Evaluation protocol}
\label{sec:protocol}

A distributional claim must specify the predicted quantity, the reference
sample and the estimator. The historical SAOI comparison uses distances and
diagnostics of scalar peak-to-valley values. The buckling scoring module
implements a fair energy score on spectral descriptors. These are different
evaluation scopes; neither a scalar statistic nor a phase-discarding descriptor
identifies the full displacement-field law.

\subsection{Proper scores and the ensemble-size correction}
\label{sec:protocol:proper}

Let $\{\mathbf{x}_i\}_{i=1}^{m}$ be the model's ensemble for a scene and
$\mathbf{y}$ an observation. Evaluating the energy score of the empirical
ensemble distribution gives
\begin{equation}
  \mathrm{ES}\bigl(\{\mathbf{x}_i\},\mathbf{y}\bigr)
  = \frac{1}{m}\sum_{i=1}^{m}\bigl\|\mathbf{x}_i-\mathbf{y}\bigr\|_2
  \;-\;\frac{1}{2m^2}\sum_{i=1}^{m}\sum_{j=1}^{m}
       \bigl\|\mathbf{x}_i-\mathbf{x}_j\bigr\|_2
  \label{eq:es}
\end{equation}
The population energy score is strictly proper for multivariate distributions
with finite first moment, and its univariate specialisation is the continuous
ranked probability score (CRPS) \citep{gneiting2007strictly}. If the ensemble
members are independent draws from a model distribution, then
\eqref{eq:es}, viewed as an estimator of that distribution's score, has
upward bias $\mathbb{E}\|X-X'\|/(2m)$. The buckling module therefore defaults
to the \emph{fair} estimator for $m\geq2$, replacing $m^2$ by $m(m-1)$
\citep{ferro2014fair}:
\begin{equation}
  \widetilde{\mathrm{ES}}
  = \frac{1}{m}\sum_{i}\bigl\|\mathbf{x}_i-\mathbf{y}\bigr\|_2
  \;-\;\frac{1}{2m(m-1)}\sum_{i}\sum_{j}\bigl\|\mathbf{x}_i-\mathbf{x}_j\bigr\|_2 .
  \label{eq:fair-es}
\end{equation}
This is unbiased over independent model draws for the population score
against a fixed observation. Fairness is not guaranteed for arbitrarily
dependent ensemble members. For a singleton deterministic forecast, the
implementation returns its distance to the observation. For independent
standard-normal forecasts and observations, the expected fair CRPS is
$1/\sqrt{\pi}$, while the expected empirical-ensemble CRPS is
$(1+1/m)/\sqrt{\pi}$. This analytic example isolates the ensemble-size bias
without relying on a particular Monte-Carlo realisation.

\subsection{Reference self-score and mean-only baseline}
\label{sec:protocol:baselines}

The energy score inherits the units and scaling of its descriptor. Even at
the true population distribution $Q$, its expected value is generally
positive: $\mathbb{E}_{Y\sim Q}\mathrm{ES}(Q,Y)=
\tfrac12\mathbb{E}_{Y,Y'\sim Q}\|Y-Y'\|$. This reflects intrinsic spread,
not merely finite reference-sample error. The buckling report provides two
empirical anchors for interpreting this value.

The quantity stored as \texttt{self\_floor} is a \emph{reference self-score}:
the implementation randomly splits the reference sample into two disjoint
halves, scores one half against every member of the other with the fair
estimator, and averages over $20$ splits (default random seed $0$). It
estimates the true-distribution score but is not a lower bound on any finite
realisation of a model's score. Repeated splits of the same sample are also
not independent replication. The \emph{mean-only reference} scores the sample's
own mean as a singleton forecast against that same sample. It is a descriptive
oracle baseline using evaluation outcomes, not a separately trained regressor.
The implemented skill transformation is
\begin{equation}
  \mathrm{SS} = \frac{\mathrm{ES}_{\mathrm{ref}}-\mathrm{ES}}
                     {\mathrm{ES}_{\mathrm{ref}}-\mathrm{ES}_{\mathrm{floor}}}
  \label{eq:skill}
\end{equation}
It equals $1$ at the estimated self-score and $0$ at the mean-only reference.
It can exceed $1$ or be negative; these anchors do not bound finite-sample
scores. A non-finite or nearly zero denominator produces an undefined skill.
Its uncertainty depends on the finite reference and generated ensembles.

\subsection{Dispersion is a separate question from location}
\label{sec:protocol:dispersion}

A score alone does not attribute error to ensemble location or width. The
current SAOI ranking script, \texttt{configs/campaigns/rank\_arms.py}, reads
saved scalar peak-to-valley samples and reports the following diagnostics.

The \emph{dispersion ratio} $\mathrm{sd}(\text{gen})/\mathrm{sd}(\text{ref})$
has population target $1$; below $1$ the generated scalar values have less
empirical spread. Both standard deviations use divisor $n$, as in
\texttt{numpy.std} with its default \texttt{ddof=0}. The \emph{bias}
$\bigl(\overline{\text{gen}}-\overline{\text{ref}}\bigr)/\mathrm{sd}(\text{ref})$
has target $0$; its magnitude measures a mean offset in reference standard
deviations. Neither diagnostic certifies distributional calibration.

The empirical rank statistic is
$r_j=m^{-1}\sum_i\mathbb{1}\{x_i<y_j\}$, using strict less-than comparison
without random tie-breaking. The script reports its distance from a uniform
CDF and the fraction with $r_j\leq0.02$ or $r_j\geq0.98$. With continuous
i.i.d.\ draws from a calibrated distribution, ranks are uniform on the
$m+1$ discrete rank values; the corresponding tail mass approaches $0.04$
for large $m$. Finite ensemble size, ties and reuse of one forecast ensemble
affect the diagnostic, so the reported KS value is descriptive, not an
independent-observation significance test. U-shaped or domed histograms can
suggest under- or over-dispersion, but bias and heterogeneous conditional
errors can also change their shape. Rank diagnostics and dispersion ratios
are complementary, not statistically independent tests.

The buckling module instead reports \texttt{spread\_ratio}, computed as
$\sqrt{\operatorname{mean}_f\operatorname{var}_{\mathrm{gen}}(X_f)}$
(sample variance with divisor $m-1$) divided by the RMS distance of reference
descriptors from the generated mean. Its denominator includes mean bias;
this is a spread-to-RMSE diagnostic, not the SAOI standard-deviation ratio.
For independent calibrated draws and $m\geq2$, the ratio of its population
second-moment scales is $\sqrt{m/(m+1)}$, which the module reports as an
i.i.d.\ moment reference. This is not the expected value of the finite-sample
ratio itself.

\subsection{Phase-discarding buckling descriptors}
\label{sec:protocol:symmetry}

The buckling energy-score implementation maps each radial-displacement field,
expressed in units of shell thickness, to circumferential Fourier magnitudes.
It takes complete even-indexed axial rows, subtracts each row's mean, retains
harmonics $1$ through $40$ by default, and concatenates the unnormalised FFT
magnitudes without per-draw amplitude normalisation. Thus it retains spectral
amplitude information while discarding the axisymmetric component, phase and
higher harmonics.

Fourier magnitudes are invariant to cyclic shifts of each sampled ring, but
are not a complete quotient by a single physical rotation: different
relative phases between harmonics or axial stations can give the same
descriptor. Nor does a deterministic imperfection in a fixed mesh establish
an exactly $SO(2)$-invariant conditional law. The score is strictly proper for
the \emph{descriptor distribution}, not for the full field or its buckle
location. Raw-field proper scoring remains valid when rotational variation
is part of the prediction target; discarding phase is an evaluation choice,
not a requirement of proper scoring. When mode labels are supplied, the
module additionally reports total variation and Jensen--Shannon divergence
between their empirical histograms.

\subsection{Scalar distances and repeated conditional observations}
\label{sec:protocol:crps-caveat}

CRPS remains proper with one observation per condition: propriety is an
expectation over verifying observations, not a requirement for repeated
outcomes at an identical condition. Limited data can nevertheless make
conditional calibration and model comparisons uncertain. Repeated
realisations at the same condition permit direct estimation of conditional
spread, which is why the historical SAOI design used $125$ outcomes per
evaluated part. These are finite reference samples, not the exact conditional
law.

The SAOI note ranks models by a normalised scalar Wasserstein distance,
not by CRPS or field-space energy score. The current ranking script
approximates the one-dimensional quantile integral at $512$ midpoint
probabilities using interpolated empirical quantiles:
\begin{equation}
 \widehat W_1/\widehat\sigma_{\mathrm{ref}}
 = \frac{1}{512\widehat\sigma_{\mathrm{ref}}}
   \sum_{k=1}^{512}\left|
    \widehat F_{\mathrm{gen}}^{-1}\!\left(\frac{k-1/2}{512}\right)
   -\widehat F_{\mathrm{ref}}^{-1}\!\left(\frac{k-1/2}{512}\right)\right|.
 \label{eq:saoi-w1}
\end{equation}
It averages per-part values to rank arms. This is an empirical distribution
distance with sampling and quadrature error; the fair-score correction in
\eqref{eq:fair-es} does not apply to it. The historical configurations and
raw score arrays are unavailable in the current checkout, so the historical
table in Section~\ref{sec:results:saoi} is attributed to the research note
rather than presented as a newly reproduced evaluation.
